% GENERATED FILE. Edit canonical lessons, then run npm run generate:books.
% canonical-source-hash: fnv1a64:6b38f75b604659d6
% canonical-lessons: PT-C518-coro, PT-C518-melodia, PT-C518-refrao, PT-C518-cantarolar, PT-C518-assobiar

\chapter{Cantar em coro}
\label{ch:pt-cantar-em-coro}
\hlchaptermodality{518}

\begin{chapteropening}
\textbf{By the end of this chapter:} \emph{I can talk about choirs, melodies and humming.}

Complete the last lesson of chapter 518: I can talk about choirs, melodies and humming.
\end{chapteropening}

\section[o coro]{o coro — choir; chorus}

\label{lesson:PT-C518-coro}



\begin{quote}
Before the new one: say the Portuguese for accordion, then the Portuguese for to tune.
\end{quote}

\subsection*{You'll want to know: o coro}
\textbf{o coro} — \textquotedblleft{}choir; chorus\textquotedblright{}.

\begin{quote}
Canto num coro há dez anos.
\end{quote}

\begin{quote}
\emph{I have sung in a choir for ten years.}
\end{quote}

\emph{Em coro} — all together.

\begin{cousinweb}
\emph{coro} is inherited through Old Portuguese (EtymDB) from Latin \emph{chorus}, from Greek \emph{khorós} \textquotedblleft{}a dance, a band of dancers and singers\textquotedblright{}.
\end{cousinweb}

\begin{grammarlens}[title={What you've built}]
One more word for telling, explaining and describing at B1.
\end{grammarlens}

\subsection*{Your turn}
\begin{itemize}
\raggedright
  \item \emph{Say it:} \emph{o coro}
  \item \emph{Say it:} \emph{o coro}, once more
  \item \emph{Say it:} a sentence of your own with \emph{o coro}
  \item \emph{Recall:} say the Portuguese for accordion, then the Portuguese for to tune, then say \emph{o coro} again
\end{itemize}

\begin{tcolorbox}[breakable,colback=teal!4,colframe=teal!35!black,title={Before you move on}]
What does o coro mean? (\textquotedblleft{}Choir; chorus\textquotedblright{}.) Say it once more.
\end{tcolorbox}
\section[a melodia]{a melodia — melody, tune}

\label{lesson:PT-C518-melodia}



\begin{quote}
Before the new one: say the Portuguese for to tune, then the Portuguese for choir.
\end{quote}

\subsection*{You'll want to know: a melodia}
\textbf{a melodia} — \textquotedblleft{}melody, tune\textquotedblright{}.

\begin{quote}
Esta melodia não me sai da cabeça.
\end{quote}

\begin{quote}
\emph{I cannot get this tune out of my head.}
\end{quote}

The adjective is \emph{melodioso}.

\begin{cousinweb}
\emph{melodia} comes from Latin \emph{melōdia} (EtymDB), from Greek \emph{melōidía} \textquotedblleft{}singing\textquotedblright{}, built on \emph{mélos} \textquotedblleft{}song\textquotedblright{}.
\end{cousinweb}

\begin{grammarlens}[title={What you've built}]
One more word for telling, explaining and describing at B1.
\end{grammarlens}

\subsection*{Your turn}
\begin{itemize}
\raggedright
  \item \emph{Say it:} \emph{a melodia}
  \item \emph{Say it:} \emph{a melodia}, once more
  \item \emph{Say it:} a sentence of your own with \emph{a melodia}
  \item \emph{Recall:} say the Portuguese for to tune, then the Portuguese for choir, then say \emph{a melodia} again
\end{itemize}

\begin{tcolorbox}[breakable,colback=teal!4,colframe=teal!35!black,title={Before you move on}]
What does a melodia mean? (\textquotedblleft{}Melody, tune\textquotedblright{}.) Say it once more.
\end{tcolorbox}
\section[o refrão]{o refrão — chorus, refrain}

\label{lesson:PT-C518-refrao}



\begin{quote}
Before the new one: say the Portuguese for choir, then the Portuguese for melody.
\end{quote}

\subsection*{You'll want to know: o refrão}
\textbf{o refrão} — \textquotedblleft{}chorus, refrain\textquotedblright{}.

\begin{quote}
Toda a gente canta o refrão.
\end{quote}

\begin{quote}
\emph{Everybody sings the chorus.}
\end{quote}

The part of a song that comes back.

\begin{cousinweb}
\emph{refrão} came through Spanish \emph{refrán} from Old French \emph{refrain} (EtymDB), built on a verb from Latin \emph{frangere} \textquotedblleft{}to break\textquotedblright{}: the song breaks off and starts again.
\end{cousinweb}

\begin{grammarlens}[title={What you've built}]
One more word for telling, explaining and describing at B1.
\end{grammarlens}

\subsection*{Your turn}
\begin{itemize}
\raggedright
  \item \emph{Say it:} \emph{o refrão}
  \item \emph{Say it:} \emph{o refrão}, once more
  \item \emph{Say it:} a sentence of your own with \emph{o refrão}
  \item \emph{Recall:} say the Portuguese for choir, then the Portuguese for melody, then say \emph{o refrão} again
\end{itemize}

\begin{tcolorbox}[breakable,colback=teal!4,colframe=teal!35!black,title={Before you move on}]
What does o refrão mean? (\textquotedblleft{}Chorus, refrain\textquotedblright{}.) Say it once more.
\end{tcolorbox}
\section[cantarolar]{cantarolar — to hum, to sing softly}

\label{lesson:PT-C518-cantarolar}



\begin{quote}
Before the new one: say the Portuguese for melody, then the Portuguese for chorus.
\end{quote}

\subsection*{You'll want to know: cantarolar}
\textbf{cantarolar} — \textquotedblleft{}to hum, to sing softly\textquotedblright{}.

\begin{quote}
A minha mãe cantarola enquanto cozinha.
\end{quote}

\begin{quote}
\emph{My mother hums while she cooks.}
\end{quote}

Built on \emph{cantar}.

\begin{cousinweb}
\emph{cantarolar} is built on \emph{cantar} \textquotedblleft{}to sing\textquotedblright{} (EtymDB), which is inherited from Latin \emph{cantāre} \textquotedblleft{}to sing\textquotedblright{}.
\end{cousinweb}

\begin{grammarlens}[title={What you've built}]
One more word for telling, explaining and describing at B1.
\end{grammarlens}

\subsection*{Your turn}
\begin{itemize}
\raggedright
  \item \emph{Say it:} \emph{cantarolar}
  \item \emph{Say it:} \emph{cantarolar}, once more
  \item \emph{Say it:} a sentence of your own with \emph{cantarolar}
  \item \emph{Recall:} say the Portuguese for melody, then the Portuguese for chorus, then say \emph{cantarolar} again
\end{itemize}

\begin{tcolorbox}[breakable,colback=teal!4,colframe=teal!35!black,title={Before you move on}]
What does cantarolar mean? (\textquotedblleft{}To hum, to sing softly\textquotedblright{}.) Say it once more.
\end{tcolorbox}
\section[assobiar]{assobiar — to whistle}

\label{lesson:PT-C518-assobiar}



\begin{quote}
Before the new one: say the Portuguese for chorus, then the Portuguese for to hum.
\end{quote}

\subsection*{You'll want to know: assobiar}
\textbf{assobiar} — \textquotedblleft{}to whistle\textquotedblright{}.

\begin{quote}
O carteiro assobia a mesma canção todas as manhãs.
\end{quote}

\begin{quote}
\emph{The postman whistles the same song every morning.}
\end{quote}

A whistle with the mouth is \emph{um assobio}.

\begin{cousinweb}
\emph{assobiar} is usually traced to Latin \emph{sībilāre} \textquotedblleft{}to hiss, to whistle\textquotedblright{} (Lewis \& Short).
\end{cousinweb}

\begin{grammarlens}[title={What you've built}]
One more word for telling, explaining and describing at B1.
\end{grammarlens}

\subsection*{Your turn}
\begin{itemize}
\raggedright
  \item \emph{Say it:} \emph{assobiar}
  \item \emph{Say it:} \emph{assobiar}, once more
  \item \emph{Say it:} a sentence of your own with \emph{assobiar}
  \item \emph{Recall:} say the Portuguese for chorus, then the Portuguese for to hum, then say \emph{assobiar} again
\end{itemize}

\begin{tcolorbox}[breakable,colback=teal!4,colframe=teal!35!black,title={Before you move on}]
What does assobiar mean? (\textquotedblleft{}To whistle\textquotedblright{}.) Say it once more.
\end{tcolorbox}
